\documentclass[a4paper,11pt]{article}

\usepackage[ngerman]{babel}
\usepackage[top=3cm,bottom=3cm,left=3cm,right=3cm]{geometry} %NICHT IN KEMIE2


%\usepackage{microtype} %schöneres typesetting  mit [stretch=10] für nur 1% anstelle von 2%
%\usepackage{lmodern} %Latin Modern FONT
%\usepackage{pifont} %schönere Font



\usepackage{amsmath}
\usepackage{nicefrac}
\usepackage{amssymb}
%\usepackage{mathtools}

\usepackage{titlesec} %Um Section, etc bearbeiten zu können
\usepackage{setspace}
\usepackage{enumitem} %enumerate
\usepackage{multicol}
\usepackage[many]{tcolorbox}
\usepackage{hyperref}
\usepackage{arydshln}

%\renewcommand{\ttdefault}{lmtt} %Schriftart wird geändert zu Latin Modern Typewriter


\hypersetup{hidelinks,
			colorlinks=true,
			linkcolor=black,
			citecolor=miudiutex,
			urlcolor=miugray2
			}

\setlength{\parindent}{0pt}
\setcounter{section}{0}


\titlespacing*{\section}{0pt}{20mm}{6mm}
\titlespacing*{\subsection}{0pt}{6mm}{3mm}
\titlespacing*{\subsubsection}{0pt}{3mm}{0mm}



%\definecolor{miudiutex}{HTML}{2f3d39}
\definecolor{miudiutex}{HTML}{1d2926}
\definecolor{miugray}{HTML}{b4cccf}
\definecolor{miugray2}{HTML}{4d7365}
\definecolor{red}{HTML}{cf1f5c}
\definecolor{green}{HTML}{00A36C}
\definecolor{delphinidin}{HTML}{315794}
\definecolor{pelargonidin}{HTML}{b33807}
\definecolor{cyanidin}{HTML}{ba0746}
\definecolor{petunidin}{HTML}{b56fed}



\raggedcolumns %wichtig für Multicolumns


%================================
% SYMBOLS
%================================

\newcommand{\RA}{$\rightarrow$~}
\newcommand{\RL}{$\rightleftharpoons$~}
\newcommand{\HARP}{\ding{226}~}
%$\xrightarrow{2}$ %Pfeil mit 2 drüber


%================================
% SHORT COMMANDS (BOXES ; ETC)
%================================


\newcommand{\notiz}[1]{\begin{notizz}{}{}\begin{center}
			{\color{miugray2!50!miudiutex}#1} \end{center}\end{notizz}}
\newcommand{\zit}[2]{\begin{zitat*}{#1}{}\small #2 \end{zitat*}}

\newcommand{\frage}[2]{\begin{qst*}{#1}{}\small #2 \end{qst*}}

\newcommand{\unten}{\end{center} \tcblower \begin{center}}


\newcommand*\circled[1]{\tikz[baseline=(char.base)]{
            \node[shape=circle,draw,inner sep=0.5pt,miudiutex] (char) {\tiny #1};}}
            
\newcommand{\miusquare}{{\color{miugray2!50}\scriptsize $\blacksquare$}~}
\newcommand{\miusquarp}{{\color{petunidin!50}\scriptsize $\blacktriangleright$}~}
\newcommand{\niu}{\item[\miusquare]}
\newcommand{\nip}{\item[\miusquarp]}

%%%%% ABSAETZE

\newcommand{\pps}{\par \vspace*{1mm}}
\newcommand{\pp}{\par \vspace*{3mm}}
\newcommand{\ppbig}{\par \vspace*{8mm}}
\newcommand{\pphuge}{\par \vspace*{10mm}}

%================================
% FRAGE BOX
%================================

\tcbuselibrary{theorems,skins,hooks}
\newtcbtheorem{qst}{Frage:}
{%
	enhanced,
	breakable,
	colback = gray!15!white,
	frame hidden,
	boxrule = 0sp,
	borderline west = {2.5pt}{0pt}{red!70},
	sharp corners,
	detach title,
	before upper = \tcbtitle\par\smallskip,
	coltitle = black,
	fonttitle = \bfseries\sffamily,
	description font = \mdseries,
	separator sign none,
	segmentation style={solid, gray!50!black},
}
{th}




%================================
% ZITAT BOX
%================================

\tcbuselibrary{theorems,skins,hooks}
\newtcbtheorem{zitat}{Zitat}
{%
	enhanced,
	breakable,
	colback = gray!15!white,
	frame hidden,
	boxrule = 0sp,
	borderline west = {2.5pt}{0pt}{black!70},
	sharp corners,
	detach title,
%	before upper = \tcbtitle\par\smallskip,
	coltitle = gray!50!black,
	fonttitle = \bfseries\sffamily,
	description font = \mdseries,
%	separator sign none,
	segmentation style={solid, gray!50!black},
}
{th}



%================================
% SIMPLE SCHWARZER RAHMEN BOX
%================================


\newcommand{\boxfr}[1]{
	\begin{tcolorbox}[
	drop shadow southeast,
	enhanced,
	colframe=black!50,
	colback=white,
	boxrule = 1pt, 
%	toprule=0pt, bottomrule=0pt,rightrule=0pt,leftrule=0pt,
	boxsep=5pt,
	arc=1pt,
	]
	\begin{center}
	#1
	\end{center}
	\end{tcolorbox}
}

\definecolor{miudiutex}{HTML}{1d2926}
\definecolor{miugray}{HTML}{b4cccf}
\definecolor{miugray2}{HTML}{4d7365}
\definecolor{white2}{HTML}{f5f7f8}
\definecolor{red}{HTML}{cf1f5c}
\definecolor{green}{HTML}{00A36C}
\definecolor{delphinidin}{HTML}{315794}
\definecolor{uniblau}{HTML}{004e9f}
\definecolor{unigelb}{HTML}{fcba00}
\definecolor{unigrau}{HTML}{909085}
\definecolor{pelargonidin}{HTML}{b33807}
\definecolor{cyanidin}{HTML}{ba0746}
\definecolor{paeonidin}{HTML}{d15ca6}
\definecolor{petunidin}{HTML}{b56fed}
\definecolor{malvidin}{HTML}{8a2d8c}

\newcommand{\metermilli}{\mskip3mu \text{mm}}
\newcommand{\cm}{\mskip3mu \text{cm}}
\newcommand{\m}{ \text{m}}
\newcommand{\lite}{ \text{l}}
\newcommand{\kg}{ \text{kg}}
\newcommand{\g}{ \text{g}}
\newcommand{\km}{ \text{km}}
\newcommand{\h}{ \text{h}}
\newcommand{\s}{ \text{s}}
\newcommand{\tonne}{ \text{t}}
\usepackage[utf8]{inputenc}
\usepackage[T1]{fontenc}
\usepackage{lmodern}

\usepackage[
    activate={true,nocompatibility},
    final,
    tracking=true,
    kerning=true,
    spacing=true,
    factor=1100,
    stretch=30,
    shrink=20
]{microtype}
\usepackage{pdfpages}
\usepackage{awesomebox}
\usepackage{marginnote}
\tcbuselibrary{documentation}
\usepackage{sidenotes}
\usepackage{import}
\usepackage{xifthen}
\usepackage{transparent}
\usepackage[table]{xcolor}


\newcommand{\abbicon}{%
  \protect\tikz[baseline=0.2mm,scale=0.8]{%
    % Das blaue Quadrat mit abgerundeten Ecken
    \fill[uniblau, rounded corners=1.2pt] (0,0) rectangle (0.8em, 0.8em);
    % Der weiße Kreis in der Mitte
    \fill[white] (0.4em, 0.4em) circle (0.12em);
  }%
} 
\usepackage[labelfont={bf},font={color=miudiutex,small},figurename={\abbicon $~$Abbildung}]{caption}


\newcommand{\bckgrnd}{
\AddToHookNext{shipout/background}{
\begin{tikzpicture}[remember picture, overlay]
 \draw[draw=none,fill=uniblau!70, shift={(current page.north west)}] (-3,0) rectangle (23,-3.75);
 \draw[draw=none,fill=miugray, shift={(current page.north east)}] (-7,0) rectangle (3,-3.75);
\end{tikzpicture}
}
}



\newcommand{\incfig}[2]{%
\def\svgwidth{#2\columnwidth}
\import{C://LaTeX-Files/pngs/}{#1.pdf_tex}
}

\renewcommand{\textbf}[1]{\textsf{{\bfseries {#1}}}}



\titleformat{\section}[hang]
		{\LARGE \color{uniblau!75!black} \sffamily \bfseries}
		{\thesection}
		{6mm}
		{}
		[]
\titleformat{\subsection}[hang]
		{\large \sffamily \bfseries \color{black}}
		{\thesubsection}
		{4mm}
		{{\color{miugray}$\blacksquare ~$}}

\titleformat{\subsubsection}[block]
		{\normalsize \bfseries \sffamily \color{miudiutex}}%
		{\normalsize \color{black} \thesubsubsection}
		{2mm}
		{{\color{miugray}$\blacksquare ~$}}
		[ \color{black}]
		
\titleformat{\paragraph}
		{\normalsize \bfseries \sffamily \color{black}}
		{\footnotesize \theparagraph}
		{1mm}
		{}

\pagestyle{empty}
\titlespacing*{\section}{0pt}{20mm}{12mm}
\titlespacing{\section}{0pt}{20mm}{12mm}

\titlespacing*{\subsection}{0pt}{14mm}{4mm}
\titlespacing{\subsection}{0pt}{14mm}{4mm}

\titlespacing*{\subsubsection}{0pt}{6mm}{3mm}
\titlespacing{\subsubsection}{0pt}{6mm}{3mm}

\titlespacing*{\paragraph}{0pt}{6mm}{2mm}
\titlespacing{\paragraph}{0pt}{6mm}{2mm}


\let\oldmarginfigure\marginfigure
\let\endoldmarginfigure\endmarginfigure

\let\oldmarginfigure\marginfigure
\let\endoldmarginfigure\endmarginfigure

\renewenvironment{marginfigure}[1][0pt] % [1] steht für 1 Argument, [0pt] ist der Standardwert
  {\oldmarginfigure[#1]\captionsetup{labelfont={sf,bf,color=uniblau!75!black},font={color=miudiutex,small}}}
  {\endoldmarginfigure}
  
  
  
  
\renewcommand{\labelitemi}{\color{uniblau!70}$\bullet$}
\renewcommand{\labelitemii}{\color{uniblau!70}$\cdot$}
\newcommand{\plmtsquare}{{\color{uniblau!70}$\blacksquare~$}}
\newcommand{\splmt}[1]{{\color{uniblau!75!black} \sffamily #1}}
\newcommand{\fplmt}[1]{{\color{uniblau!75!black} \sffamily \bfseries #1}}

\newcommand{\antwort}[1]{
\begin{tcolorbox}[
	enhanced,
	enforce breakable,
	standard jigsaw,
	boxrule=0.7pt,
	colframe=black,
	sharp corners,
	boxsep=5pt,
	title=\bfseries \sffamily Antwort,
	opacityback=0,
	coltitle=miudiutex,
	attach boxed title to top text left={yshift=-\tcboxedtitleheight/2},
	boxed title style={colback=white,colframe=white},
	bottomrule at break=0pt,
	toprule at break=0pt,
	pad at break=16mm,
]
	\begin{center}
	#1
	\end{center}
	\end{tcolorbox}
}




\rowcolors{2}{miugray!50}{miugray!50}
\arrayrulecolor{white} % Weiße Gitterlinien
\setlength{\arrayrulewidth}{1.5pt} % Etwas dickere Linien
\renewcommand{\arraystretch}{1.4} % Mehr Zeilenhöhe für bessere Lesbarkeit}


%\setlist[itemize]{itemsep=0mm}
\setlist[enumerate]{itemsep=0mm}
\usepackage[table]{xcolor}
\usepackage{booktabs}
\usepackage{pgfplots}
    \pgfplotsset{
        every axis post/.style={
	yticklabel=\empty,
	xticklabel=\empty,
	y label style={at={(axis description cs:-0.2,1.1)}},
	x label style={at={(axis description cs:1.05,0.2)}},
	ticks=none,
	width=6cm,
	axis lines=middle,
%	xlabel near ticks,
        },
    }
\usetikzlibrary{positioning, arrows.meta, calc, patterns}

\usepackage{relsize}
\setlength{\abovedisplayskip}{8pt}
\setlength{\belowdisplayskip}{8pt}
\setlength{\abovedisplayshortskip}{8pt}
\setlength{\belowdisplayshortskip}{8pt}

\begin{document}
\newgeometry{top=1cm,bottom=3cm,left=4cm,right=4cm}

\AddToHookNext{shipout/background}{
\begin{tikzpicture}[remember picture, overlay]
 \draw[draw=none,fill=uniblau!70, shift={(current page.north west)}] (-3,0) rectangle (26,-2.75);
% \draw[draw=none,fill=miugray, shift={(current page.north east)}] (-1,0) rectangle (3,-3.75);
\end{tikzpicture}
}
\begin{center}
\LARGE
\color{white}
\textbf{Theorie} \par \vspace*{2mm}
\large \color{miugray} \textbf{Prozessbezogene Lebensmitteltechnologie}
\end{center}
\par 
\vspace*{4mm}


\color{miudiutex}

\section*{Sinkgeschwindigkeit nach Stokes}
\subsection*{Grundlagen}
Für ein kugelförmiges Teilchen im Fluid gelten folgende geometrische Beziehungen für das Volumen $V_{\text{Kugel}}$ und die Fläche $A$:
\begin{align*}
    V_{\text{Kugel}} \qquad &= \frac{4}{3}\pi r^3 \qquad = \frac{1}{6}\pi d^3 \\
    A \qquad &= \pi r^2 \qquad ~~  = \frac{1}{4}\pi d^2
\end{align*}

\subsection*{Wirksame Kräfte}

    \textbf{Gewichtskraft $F_G$:} Über die Masse des Teilchens $m_T$ und dessen Dichte $\rho_T$:
    \begin{equation}
        F_G = m_T \cdot g = \frac{\pi}{6} \cdot d^3 \cdot \rho_T \cdot g
    \end{equation}
    \textbf{Auftriebskraft $F_A$:} Entspricht der Gewichtskraft des verdrängten Fluidvolumens (Fluiddichte $\rho_F$):
    \begin{equation}
        F_A = m_F \cdot g = \frac{\pi}{6} \cdot d^3 \cdot \rho_F \cdot g
    \end{equation}
   \textbf{Widerstandskraft $F_W$:} Abhängig vom Strömungswiderstandsbeiwert $c_W$ und der Sinkgeschwindigkeit $v$:
    \begin{equation}
        F_W = c_W \cdot \frac{\rho_F}{2} \cdot v^2 \cdot A = c_W \cdot \frac{\rho_F}{2} \cdot v^2 \cdot \frac{\pi}{4} d^2
    \end{equation}
\textbf{Trägheitskraft $F_P$:} Beschleunigungskraft inklusive der hydrodynamisch ''mitbewegten'' Fluidmasse:
    \begin{equation}
        F_P = \left(m_T + \frac{m_F}{2}\right) \cdot \frac{dv}{dt} = \frac{\pi}{6} \cdot d^3 \cdot \left(\rho_T + \frac{\rho_F}{2}\right) \cdot \frac{dv}{dt}
    \end{equation}

\pagebreak
\newgeometry{top=3cm,bottom=3cm,left=4cm,right=4cm}


\subsection*{Welche Kräfte wirken auf das Teilchen ein?}
Im stationären Zustand (konstante Sinkgeschwindigkeit, $F_P = 0$) hält die Gewichtskraft den entgegengesetzten Kräften (Auftrieb und Widerstand) das Gleichgewicht:
\begin{equation}
    F_G = F_A + F_W
\end{equation}
\begin{figure}[h]
\centering
\begin{tikzpicture}[scale=1.6, >=stealth]
    % Wellenlinien links und rechts (Fluid-Andeutung)
%    \draw[domain=-1:1, smooth, variable=\y, shift={(-1.2,0)}] plot ({0.1*sin(\y*400)}, \y);
%    \draw[domain=-1:1, smooth, variable=\y, shift={(1.2,0)}] plot ({0.1*sin(\y*400)}, \y);
    \draw[fill=uniblau!10,draw=miugray!30] (-3,1.65) rectangle (3,-1.65);
    % Kugel
    \draw[fill=miugray!80, draw=black, thick] (0,0) circle (0.5cm);
    \node[scale=0.5] at (0.1,-0.25) {$m_T, \rho_T$};
    
    % Kraftpfeile nach oben
    \draw[->, thick] (0, 0.5) -- (0, 1.1) node[above] {$F_W$};
    \draw[->, thick] (0.3, 0.4) -- (0.3, 0.8) node[above] {$F_A$};
    \draw[->, thick] (-0.3, 0.4) -- (-0.3, 0.8) node[above] {$F_P$};
    
    % Kraftpfeil nach unten
    \draw[->, thick] (0, -0.5) -- (0, -1.1) node[below] {$F_G$};
    
    % Sinkgeschwindigkeit Pfeil
    \draw[->, dashed, black, thick] (0.3, -0.4) -- (0.3, -0.8) node[right] {$v_{\text{sinkgeschw.}}$};
    
	%Durchmesser Kugel Pfeil
	 \draw[<->, dotted, black, thick,rotate=35]   (-0.5,0) -- (0.5,0) node[midway,yshift=2.5mm,xshift=-1mm,scale=0.8]{$d$};
    
    % Fluid Parameter
    \node[left] at (-1.5, -1) {$\eta_F, \rho_F$};
\end{tikzpicture}
\caption{Schematische Darstellung eines kugelförmigen Teilchen in einem Fluid. \textit{Hinweis:} Bei kleinen Teilchen mit einem Durchmesser $d < 1\,\text{mm}$ kann die Trägheitskraft $F_P$ für den stationären Zustand vernachlässigt werden ($d v/dt \approx 0$).}
\end{figure}



%\section{Herleitung Sinkgeschwindigkeit}


\subsection*{Wie groß ist also die Sinkgeschwindigkeit $v$?}
Durch Einsetzen der Definitionen in das Kräftegleichgewicht erhalten wir:
\begin{align*}
    F_G - F_A &= F_W \\
    \frac{\pi}{6} d^3 \rho_T g - \frac{\pi}{6} d^3 \rho_F g &= c_W \cdot \frac{\rho_F}{2} \cdot v^2 \cdot \frac{\pi}{4} d^2 \\
    \frac{\pi}{6} d^3 g (\rho_T - \rho_F) &= c_W \cdot \frac{\pi}{8} \rho_F d^2 v^2
\end{align*}
Kürzen von $\pi$ und $d^2$ sowie Multiplikation mit 8 isoliert $v^2$:
\begin{align*}
    \frac{8}{6} d g (\rho_T - \rho_F) &= c_W \rho_F v^2 \\
    v^2 &= \frac{4 d g (\rho_T - \rho_F)}{3 \cdot c_W \cdot \rho_F}
\end{align*}
Daraus folgt die allgemeine Formel für die Sinkgeschwindigkeit:
\antwort{
\begin{equation}
    v = \sqrt{\frac{4 \cdot d \cdot g \cdot (\rho_T - \rho_F)}{3 \cdot c_W \cdot \rho_F}}
\end{equation}
}
\pagebreak
\subsection*{Abhängigkeit des Widerstandsbeiwerts $c_W$ \& Laminarer Fall}
Der Widerstandsbeiwert $c_W$ ist eine Funktion der Reynolds-Zahl ($\text{Re}$), welche das Verhältnis von Trägheits- zu Zähigkeitskräften beschreibt. Er hängt von der Körperform, -lage, Dichte, der dynamischen Viskosität $\eta_F$ und der Strömungsgeschwindigkeit ab:
\begin{align*}
    c_W &= \frac{24}{\text{Re}} \quad &&\text{für } \text{Re} < 0{,}5 \quad \text{(laminare Strömung)} \\
    c_W &= \frac{18{,}5}{\text{Re}^{0{,}6}} \quad &&\text{für } 0{,}5 < \text{Re} < 500 \\
    c_W &= 0{,}44 \quad &&\text{für } \text{Re} > 500 \quad \text{(turbulente Strömung)}
\end{align*}
\ppbig
Die Reynolds-Zahl ist definiert als:
\begin{equation}
    \text{Re} = \frac{\rho_F \cdot d \cdot v}{\eta_F}
\end{equation}

\subsubsection*{Frage: Wenn die Strömung laminar ist, wie verändert sich $v$?}
Für den laminaren Fall ($\text{Re} < 0{,}5$) setzen wir $c_W = \frac{24}{\text{Re}}$ in das Ergebnis aus Teil (b) ein:
\begin{align*}
    v^2 &= \frac{4 d g (\rho_T - \rho_F)}{3 \cdot \left(\dfrac{24 \cdot \eta_F}{\rho_F \cdot d \cdot v}\right) \cdot \rho_F} \\
    &= \frac{4 d^2 g v (\rho_T - \rho_F)}{72 \cdot \eta_F}
\end{align*}
\ppbig
Durch Teilen beider Seiten durch $v$ und Kürzen des Bruchs ($\frac{4}{72} = \frac{1}{18}$) ergibt sich das finale \textbf{Gesetz von Stokes}:
\antwort{
\begin{equation}
    v = \frac{d^2 \cdot g \cdot (\rho_T - \rho_F)}{18 \eta_F}
\end{equation}
}
\end{document}